\section{Lecture 3: Hypothesis Spaces and Inductive Bias}

\subsection{Motivation}

In the previous lectures, we formulated learning as:
\begin{itemize}
    \item approximating an unknown function,
    \item minimizing expected loss under a data distribution.
\end{itemize}

\medskip

\noindent
However, a fundamental question remains:

\medskip

\noindent
\textit{Among all possible functions, which ones should we consider?}

\medskip

\noindent
Since we only observe finite data, many functions can fit the data equally well. Without further assumptions, learning is ill-posed.

\medskip

\noindent
\textbf{Guiding idea.}  
Learning requires restricting the space of possible functions and introducing assumptions about the structure of the problem.

\subsection{Function Classes and Model Families}

We define a \emph{hypothesis space}:
\[
\mathcal{H} = \{ f_\theta : \mathcal{X} \to \mathcal{Y} \mid \theta \in \Theta \}.
\]

\medskip

\noindent
\textbf{Examples.}

\begin{itemize}
    \item \textbf{Linear models:}
    \[
    f(x) = w^\top x
    \]

    \item \textbf{Polynomial models:}
    \[
    f(x) = \sum_{k=0}^d a_k x^k
    \]

    \item \textbf{Neural networks:}
    \[
    f(x) = \sigma(W_2 \sigma(W_1 x))
    \]
\end{itemize}

\medskip

\noindent
\textbf{Interpretation.}
\begin{itemize}
    \item $\mathcal{H}$ defines which functions are representable
    \item Learning selects a function within $\mathcal{H}$
\end{itemize}

\medskip

\noindent
\textbf{Expressivity.}  
A richer hypothesis space can approximate more complex functions.

\medskip

\noindent
\textbf{Tradeoff.}
\begin{itemize}
    \item Small $\mathcal{H}$: limited expressivity
    \item Large $\mathcal{H}$: risk of overfitting
\end{itemize}

\subsection{Inductive Bias and Prior Structure}

The restriction to a hypothesis space introduces an \emph{inductive bias}.

\medskip

\noindent
\textbf{Definition.}  
Inductive bias is the set of assumptions that a learning algorithm uses to generalize beyond observed data.

\medskip

\noindent
\textbf{Sources of inductive bias.}
\begin{itemize}
    \item Choice of hypothesis space
    \item Regularization
    \item Optimization algorithm
\end{itemize}

\medskip

\noindent
\textbf{Example.}
\begin{itemize}
    \item Linear models assume linear relationships
    \item Convolutional networks assume translation invariance
\end{itemize}

\medskip

\noindent
\textbf{Prior structure.}  
Inductive bias can be interpreted probabilistically as a prior over functions:
\[
p(f).
\]

\medskip

\noindent
\textbf{Insight.}  
Learning is not purely data-driven; it depends critically on assumptions encoded in the model.

\subsection{Bias--Variance Tradeoff (Intuition)}

A central concept in learning is the tradeoff between bias and variance.

\medskip

\noindent
\textbf{Bias.}
\begin{itemize}
    \item Error due to model assumptions
    \item High when the model is too simple
\end{itemize}

\medskip

\noindent
\textbf{Variance.}
\begin{itemize}
    \item Sensitivity to data fluctuations
    \item High when the model is too complex
\end{itemize}

\medskip

\noindent
\textbf{Qualitative decomposition.}
\[
\text{Error} = \text{Bias}^2 + \text{Variance} + \text{Noise}.
\]

\medskip

\noindent
\textbf{Interpretation.}
\begin{itemize}
    \item Increasing model complexity reduces bias
    \item Increasing model complexity increases variance
\end{itemize}

\medskip

\noindent
\textbf{Example.}
\begin{itemize}
    \item Linear model on nonlinear data: high bias
    \item Highly flexible model on small dataset: high variance
\end{itemize}

\medskip

\noindent
\textbf{Insight.}  
Good generalization requires balancing simplicity and flexibility.

\subsection{Choosing Representations}

The representation of data plays a central role in learning.

\medskip

\noindent
\textbf{Feature representation.}
\begin{itemize}
    \item Input $x$ may be transformed into features $\phi(x)$
    \item Learning is performed on $\phi(x)$ rather than $x$
\end{itemize}

\medskip

\noindent
\textbf{Examples.}
\begin{itemize}
    \item Polynomial features
    \item Kernel mappings
    \item Learned representations (neural networks)
\end{itemize}

\medskip

\noindent
\textbf{Key idea.}  
The difficulty of learning depends on the representation.

\medskip

\noindent
\textbf{Linearization principle.}  
A complex problem may become simple under the right representation.

\medskip

\noindent
\textbf{Modern perspective.}
\begin{itemize}
    \item Deep learning automatically learns representations
    \item Representation learning replaces manual feature design
\end{itemize}

\medskip

\noindent
\textbf{Insight.}  
Choosing or learning representations is often more important than the choice of model.

\subsection{A Unifying Perspective}

Hypothesis spaces and inductive bias define the structure of learning:

\begin{itemize}
    \item \textbf{Hypothesis space:} what functions are possible
    \item \textbf{Inductive bias:} which functions are preferred
    \item \textbf{Data:} which functions are consistent with observations
\end{itemize}

\medskip

\noindent
\textbf{Core idea.}  
Learning is the process of selecting a function that balances data fit with prior assumptions.

\medskip

\noindent
\textbf{Key insight.}  
Generalization is possible only because of inductive bias; without it, learning from finite data would be impossible.

\subsection{Outlook}

In this lecture, we examined:

\begin{itemize}
    \item hypothesis spaces and model families,
    \item inductive bias and prior structure,
    \item the bias--variance tradeoff,
    \item the role of representations.
\end{itemize}

\medskip

\noindent
These ideas naturally lead to the next question:

\medskip

\noindent
\textit{How do we actually find the best function within a hypothesis space?}

\medskip

\noindent
In the next lecture, we study learning as an optimization problem, where parameter estimation is framed as minimizing an objective function.

\medskip

\noindent
\textbf{Guiding principle.}  
The success of learning depends as much on the choice of representation and inductive bias as on the data itself.